% !TEX root = ../main.tex

\begin{innovations}

本文围绕人形机器人的遥操作、全身运动控制与感知闭环任务执行展开研究，针对“数据难获取、控制难统一、系统难闭环”的关键问题，形成了以下几方面创新性工作。

第一，构建了面向人形机器人全身控制的多源异构人体运动数据采集与统一重定向框架。本文综合惯性动捕、光学动捕和仿真交互三类数据来源，兼顾部署灵活性、采集精度和任务多样性，并采用基于 GMR 的动作重定向方法处理人体与机器人之间在肢体尺度、关节结构和自由度分布上的差异，为后续控制策略学习与真机执行提供统一的数据接口。

第二，提出了面向移动操作场景的上下肢解耦式通用全身运动控制架构。与完全端到端的全身控制方法相比，该架构将下肢运动稳定性问题与上肢操作精度问题分而治之：下肢利用强化学习获得鲁棒的基础运动能力，上肢通过轨迹接口、逆运动学和补偿机制实现高保真动作跟踪，同时在系统层面设计了上下肢协调机制，从而兼顾控制性能、训练效率与工程可部署性。

第三，验证了从遥操作数据采集到感知--控制闭环任务执行的完整技术链路。本文不仅关注离线策略训练或单一子模块实现，而且将所提出的方法应用于 CyboRacket 动态球拍任务和 OpenClaw 语音服务任务中，实现了感知、规划与全身控制的系统级集成。这说明所提出的框架能够从示教和数据采集阶段平滑过渡到自主执行阶段，具有较好的通用性与扩展潜力。

\end{innovations}


% \begin{innovations}
%
% 工程博士学位论文应主要聚焦工程实践和应用研究，须体现工程性、创新性、实践性、应用性特征，
% 体现学位申请人在专业领域掌握坚实全面的基础理论和系统深入的专门知识，
% 具有独立承担专业实践工作的能力，在专业实践领域做出创新性成果，
% 对推动本专业领域知识和技术的发展作出重要贡献。
%
% 故在此须说明本学位论文的创新性及应用性，确保符合工程博士学位论文要求，字数在800字以内。
%
% \end{innovations}
